% GENERATED FILE. Edit canonical lessons, then run npm run generate:books.
% canonical-source-hash: fnv1a64:38eab0c81f1caed5
% canonical-lessons: GE-C255-du-solltest, GE-C255-erklaeren, GE-C255-review

\chapter{You Should — du solltest, and Explaining Why}
\label{ch:ge-you-should}
\hlchaptermodality{255}

\begin{chapteropening}
\textbf{By the end of this chapter:} \emph{I can give advice with du solltest and man sollte, and explain why in three steps.}

Complete GE-C255-review: a piece of advice with du solltest, and an explanation with warum, weil and darum.
\end{chapteropening}

\section[du solltest / man sollte]{du solltest — you should}

\label{lesson:GE-C255-du-solltest}



\begin{quote}
Say \textquotedblleft{}I can take the umbrella\textquotedblright{}. (\emph{Ich kann den Regenschirm mitnehmen.}) The modal second, the infinitive at the end.
\end{quote}

\subsection*{You'll want to know: du solltest}
\textbf{du solltest} — \textquotedblleft{}you should\textquotedblright{}: \emph{doo ZOLL-test}.

\begin{quote}
\emph{Du solltest den Regenschirm mitnehmen.} — You should take the umbrella.
\end{quote}

It is a modal, so it builds the same bracket as \emph{kann}: \emph{solltest} second, the infinitive last. Its forms: \emph{ich sollte}, \emph{du solltest}, \emph{er sollte}, \emph{wir sollten}, \emph{Sie sollten}.

\emph{Solltest} looks like a past, and it is spelled exactly like the Präteritum of \emph{sollen}. Here it is the \textbf{Konjunktiv II}, the mood of \textquotedblleft{}it would be good if\textquotedblright{}. That is what makes it advice. The plain present, \emph{du sollst}, is stronger: it passes on an order or a duty, \textquotedblleft{}you are to\textquotedblright{}.

\begin{grammarlens}[title={Grammar Lens: man sollte, and two verbs at the end}]
For advice to everybody, German uses \textbf{man}, \textquotedblleft{}one, people\textquotedblright{}:

\begin{quote}
\emph{Man sollte den Fahrplan lesen.} — One should read the timetable.
\end{quote}

\emph{Man} takes the \emph{er} form of the verb: \emph{man sollte}, like \emph{er sollte}.

After \emph{dass} or \emph{weil}, both verbs go to the end, and the modal comes last of all:

\begin{quote}
\emph{Ich finde, dass du den Regenschirm mitnehmen} \textbf{solltest}.
\end{quote}
\end{grammarlens}

\begin{cousinweb}
\textbf{sollen} $\leftarrow$ Old High German \emph{sculan} $\leftarrow$ Proto-Germanic \emph{*skulaną}, \textquotedblleft{}to owe\textquotedblright{} (EtymDB 2.1). EtymDB sets English \emph{shall} and \emph{should} beside it: \emph{should} is the past of \emph{shall}, as \emph{solltest} is of \emph{sollen}. You met the same root in \emph{Entschuldigung}, inside \emph{Schuld}, a debt or a fault: what you should do is what you owe.
\end{cousinweb}

\subsection*{Your turn}
Say these aloud:

\begin{itemize}
\raggedright
  \item \emph{du solltest}
  \item \emph{Du solltest den Regenschirm mitnehmen.}
  \item \emph{Man sollte den Fahrplan lesen.}
  \item \emph{ich sollte, du solltest, er sollte, wir sollten}
  \item \emph{Ich finde, dass du den Regenschirm mitnehmen solltest.}
\end{itemize}

\begin{tcolorbox}[breakable,colback=teal!4,colframe=teal!35!black,title={Before you move on}]
\textquotedblleft{}You should take the umbrella\textquotedblright{}? (\emph{\textbf{Du solltest den Regenschirm mitnehmen.}}) Advice, or an order: \emph{du solltest} or \emph{du sollst}? (\emph{\textbf{Du solltest}} is advice.) After \emph{dass}, which verb comes last? (\textbf{The modal}, \emph{solltest}.)
\end{tcolorbox}
\section[erklären, warum …]{erklären, warum — explaining why}

\label{lesson:GE-C255-erklaeren}



\begin{quote}
Say \textquotedblleft{}to explain\textquotedblright{}. (\emph{Erklären.}) Then ask \textquotedblleft{}why is the train late?\textquotedblright{} (\emph{Warum hat der Zug Verspätung?})
\end{quote}

\begin{grammarlens}[title={Grammar Lens: a question inside a sentence}]
Put that question inside a request and its verb moves to the end:

\begin{quote}
\emph{Warum} \textbf{hat} \emph{der Zug Verspätung?} — Why is the train late? \emph{Können Sie mir erklären, warum der Zug Verspätung} \textbf{hat}? — Can you explain to me why the train is late?
\end{quote}

Inside a sentence, \textbf{warum} opens a clause the way \emph{weil} and \emph{dass} do. \emph{Wie}, \emph{wo} and \emph{wann} behave the same way: \emph{Ich erkläre dir, wie das} \textbf{geht}. \emph{Weißt du, wann der Zug} \textbf{kommt}?

\emph{Erklären} takes the person in the dative: \emph{mir}, \emph{dir}, \emph{Ihnen}, \textquotedblleft{}to me, to you\textquotedblright{}.
\end{grammarlens}

\begin{culture}
A clear explanation in German moves the way it does in English: what, why, and what follows.

\begin{quote}
\emph{Ich komme heute nicht.} — what. \emph{Ich komme nicht, weil ich Fieber} \textbf{habe}. — why, with \emph{weil}. \emph{Darum} \textbf{rufe} \emph{ich dich an.} — what follows, with \emph{darum}.
\end{quote}

Each step has its own word order: plain, verb last, verb second. Said in a row, the three hold together as one explanation.
\end{culture}

\begin{cousinweb}
\textbf{erklären} is \emph{er-} + \emph{klar}, \textquotedblleft{}clear\textquotedblright{}: to make clear. EtymDB 2.1 traces \emph{klar} through Middle High German \emph{klār} to Latin \emph{clārus}, \textquotedblleft{}bright, clear\textquotedblright{}. The noun \textbf{die Erklärung}, the explanation, is built on the same verb.
\end{cousinweb}

\subsection*{Your turn}
Say these aloud:

\begin{itemize}
\raggedright
  \item \emph{Können Sie mir erklären, warum der Zug Verspätung hat?}
  \item \emph{Ich erkläre dir, wie das geht.}
  \item \emph{Weißt du, wann der Zug kommt?}
  \item the three steps: \emph{Ich komme heute nicht. Ich komme nicht, weil ich Fieber habe. Darum rufe ich dich an.}
  \item advice to finish it: \emph{Du solltest Fieber messen.}
\end{itemize}

\begin{tcolorbox}[breakable,colback=teal!4,colframe=teal!35!black,title={Before you move on}]
Inside \textquotedblleft{}explain to me why …\textquotedblright{}, where does the verb go? (\textbf{To the end.}) The three steps of an explanation? (\textbf{What, why with} \emph{\textbf{weil}}, \textbf{and what follows with} \emph{\textbf{darum}}.) What does \emph{erklären} make a thing? (\emph{\textbf{Klar}}, clear.)
\end{tcolorbox}
\section[(review)]{Review — what you should do, and why}

\label{lesson:GE-C255-review}



\begin{quote}
Say \textquotedblleft{}you should read the timetable\textquotedblright{}. (\emph{Du solltest den Fahrplan lesen.})
\end{quote}

\begin{grammarlens}[title={What you've built}]
Two chapters, and the whole of giving reasons is in place. An opinion, with \emph{dass} or \emph{meiner Meinung nach}. A reason with \emph{weil}, a consequence with \emph{darum}. Advice with \emph{du solltest} or \emph{man sollte}. And a question folded into a sentence, with its verb at the end: \emph{erklären, warum …}

Here is a friend's problem, and an answer that uses three of those pieces:

\begin{quote}
\emph{Der Zug hat Verspätung. Was soll ich machen?} \emph{Meiner Meinung nach solltest du ein Taxi nehmen, weil es regnet.}
\end{quote}

Listen to the answer: after \emph{meiner Meinung nach}, the modal comes straight next, \emph{solltest du}, and the infinitive \emph{nehmen} waits at the end of its clause.
\end{grammarlens}

\subsection*{Your turn}
Say these aloud:

\begin{itemize}
\raggedright
  \item \emph{Man sollte den Fahrplan lesen.}
  \item \emph{Können Sie mir erklären, warum der Zug Verspätung hat?}
  \item \emph{Meiner Meinung nach solltest du ein Taxi nehmen, weil es regnet.}
  \item \emph{Ich finde, dass du ein Taxi nehmen solltest.}
  \item your own: a piece of advice, a reason with \emph{weil}, and what follows with \emph{darum}
\end{itemize}

\begin{tcolorbox}[breakable,colback=teal!4,colframe=teal!35!black,title={Before you move on}]
\textquotedblleft{}One should\textquotedblright{}? (\emph{\textbf{Man sollte.}}) After \emph{meiner Meinung nach}, which comes next? (\textbf{The verb}, \emph{solltest du}.) Inside \emph{erklären, warum}, where is the verb? (\textbf{At the end.})
\end{tcolorbox}
